\clearpage
\phantomsection
\chapter{SYSTEM ANALYSIS AND DESIGN}

\section{Requirements Analysis}

\subsection{Functional Requirements}

The system must fulfill the following functional requirements:
\begin{enumerate}
    \item \textbf{FR1 -- Game Engine}: Implement a complete, self-contained Roguelike Deckbuilder game engine in C\# supporting: turn-based combat with a card-based attack system, a procedurally generated room map, deck acquisition and management, relic collection, multiple room types (Monster, Elite, Event, Shop, Rest), and a data-driven parameter system stored in JSON.
    \item \textbf{FR2 -- Deterministic Baseline Agent}: Provide a rule-based AI agent that can play the game to completion, serving as a playability validator and initial balancing oracle.
    \item \textbf{FR3 -- RL Environment}: Expose the game engine as an OpenAI Gymnasium-compatible environment with a well-defined observation space, action space, and action masking mechanism.
    \item \textbf{FR4 -- RL Agent Training}: Train four MaskablePPO agents with distinct reward functions and produce checkpoints that can be loaded for evaluation.
    \item \textbf{FR5 -- Game Parameter Optimization}: Implement both a single-objective GA and a multi-objective NSGA-II that use agent evaluation to search for balanced parameter configurations.
    \item \textbf{FR6 -- Results Visualization}: Generate evolution charts, Pareto front visualizations, and card viability analyses.
\end{enumerate}

\subsection{Non-Functional Requirements}

\begin{enumerate}
    \item \textbf{NFR1 -- Cross-platform Compatibility}: The game engine (C\#) must run on Linux to be usable inside Kaggle kernels for GPU training.
    \item \textbf{NFR2 -- Simulation Speed}: The game engine must support at least 1,000 complete game runs per second to make population-based optimization tractable.
    \item \textbf{NFR3 -- Reproducibility}: Experiments must be reproducible via fixed random seeds.
    \item \textbf{NFR4 -- Extensibility}: New cards, relics, and enemy types must be addable by modifying only the JSON data files without changing source code.
\end{enumerate}

\section{System Architecture}

The system is organized into three principal layers, as illustrated in Figure~\ref{fig:architecture}:

\begin{figure}[H]
    \centering
    \includegraphics[width=0.92\textwidth]{SystemStructure.png}
    \caption{High-level system architecture showing data flow between the Optimization Layer, Simulation/Training Layer, Data Layer, and Core Game Layer.}
    \label{fig:architecture}
\end{figure}

\begin{itemize}
    \item \textbf{Core Game Layer}: The C\# game engine implementing all gameplay mechanics. This layer is deterministic given a fixed random seed and a fixed parameter set $\theta$.
    \item \textbf{Simulation / Training Layer}: A Python bridge that (a) exposes the game as a Gymnasium RL environment for agent training, and (b) runs mass simulations to compute fitness metrics for the optimization layer.
    \item \textbf{Optimization Layer}: Python implementations of GA and NSGA-II that produce candidate parameter vectors $\theta$, dispatch them to the simulation layer, and evolve populations based on returned fitness values.
    \item \textbf{Data Layer}: JSON files that define all game entities (cards, enemies, relics, effects) and store optimization results. Separating data from logic enables the optimization layer to modify parameters without recompiling any code.
\end{itemize}

\section{Core Game Engine Design}

\subsection{Game Structure Overview}

A single game run consists of 15 floors organized into three Acts of five floors each. The player navigates a branching map of rooms, engaging in combat, collecting cards and relics, and managing resources (HP, gold, mana). The game ends in victory if the player defeats the final boss on Floor 15, or in defeat if HP drops to zero.

\textbf{Map generation} uses a constrained random algorithm that guarantees at least one viable path from the entrance to the exit floor while distributing room types according to configurable weights. A typical 15-floor map contains: 8--10 Monster rooms, 2--3 Elite rooms, 1--2 Shop rooms, 2--3 Rest Sites, and 2--3 Event rooms.

\textbf{Card system}: The game contains over 30 distinct cards organized by type (Attack, Skill, Power) and star rating (1--5). Cards have a mana cost, and each round the player draws five cards, spending mana to play them. Cards are acquired as rewards after combat and can be removed or upgraded at rest sites or via events.

\textbf{Relic system}: Relics are passive items collected from Elite rooms, events, and shops. Over 20 relics are defined, each with a trigger condition (e.g., \texttt{OnEnemyAttack}, \texttt{OnTurnStart}) and an effect (e.g., add thorns, reduce mana cost, draw extra cards).

\subsection{Combat System Design}

The combat system follows a turn-based structure with alternating player and enemy phases:

\begin{enumerate}
    \item \textbf{Player turn}: The player draws 5 cards, then plays any combination of affordable cards using available mana (default: 3 mana per turn). Actions include dealing damage to a target enemy, gaining block (temporary defense), applying status effects, and drawing additional cards.
    \item \textbf{Enemy turn}: Each enemy executes its scheduled action (attack, buff, or debuff) according to a pre-defined pattern. Status effects (Strength, Weakness, Poison, Burn, Vulnerability) are resolved.
    \item \textbf{End-of-turn}: Block is cleared, effects are ticked down, and the player draws a new hand.
\end{enumerate}

Combat resolves when all enemies are defeated (player wins the room) or the player's HP reaches zero (run ends). The class diagram in Figure~\ref{fig:class_diagram} shows the relationship between the key combat classes.

\begin{figure}[H]
    \centering
    \includegraphics[width=0.65\textwidth]{class_diagram.png}
    \caption{Class diagram of the core game module showing the relationships between \texttt{GameController}, \texttt{CombatManager}, \texttt{DeckManager}, the AI interface, and the optimization layer.}
    \label{fig:class_diagram}
\end{figure}

Figure~\ref{fig:combat_seq} illustrates the sequence of interactions during a single combat turn.

\begin{figure}[H]
    \centering
    \includegraphics[width=0.65\textwidth]{combat_sequence.png}
    \caption{Sequence diagram of one combat turn: the Combat Manager draws cards from the Deck Manager, queries the AI Agent for a decision, resolves the card action, executes the enemy turn, and checks win/loss conditions.}
    \label{fig:combat_seq}
\end{figure}

\section{Deterministic Heuristic Agent Design}

\subsection{Purpose and Role}

The \textbf{Deterministic Heuristic Agent (DHA)}, implemented as \texttt{HeuristicPlayerAI}, serves two specific roles in the project:
\begin{enumerate}
    \item \textbf{Playability validation}: Before any optimization or RL training, the DHA confirms that the game is winnable and that no trivially broken configurations exist. A game that a reasonable agent cannot win at all (or can always win trivially) is not ready for balancing.
    \item \textbf{Optimization benchmark}: In the initial GA/NSGA-II experiments documented in Section~\ref{sec:results_opt}, the DHA provides a consistent, fast, and reproducible evaluation signal. Because the DHA executes in microseconds per decision (no neural network inference), it enables the optimization algorithm to run orders of magnitude more simulations than would be possible with RL agents.
\end{enumerate}

\subsection{Decision Logic}

The DHA implements a three-level priority hierarchy for in-combat decisions:

\begin{enumerate}
    \item \textbf{Survival first}: Calculate total incoming damage from all enemies in the current turn. If incoming damage exceeds current HP, preferentially play defensive cards (Block, Shielding skills) to survive.
    \item \textbf{Lethal second}: If any single target can be killed this turn with available cards and mana, prioritize those cards to reduce enemy damage output.
    \item \textbf{Efficient offense third}: Otherwise, play the card with the highest damage-per-mana ratio, preferring multi-target attacks when multiple enemies are present.
\end{enumerate}

For map navigation, the DHA scores each available next-room type based on current HP percentage: it seeks Elite rooms when HP $\geq$ 60\% (rewarded with rare cards and relics), Rest sites when HP $\leq$ 40\%, and Monster rooms otherwise. Card rewards are chosen by star rating, preferring higher rarities and maintaining a balanced ratio of attack to skill cards.

The DHA's deterministic nature makes it an ideal baseline: given the same game seed and parameters, it always makes the same choices, ensuring that measured fitness differences between parameter sets are attributable to the parameters rather than agent stochasticity.

\section{Reinforcement Learning Environment Design}

\subsection{Observation Space}

The observation vector $\mathbf{s} \in \mathbb{R}^{148}$ encodes the full game state at the moment the agent must make a decision. Table~\ref{tab:obs_space} describes its components.

\begin{table}[H]
    \centering
    \caption{Observation space components (total: 148 dimensions).}
    \label{tab:obs_space}
    \begin{tabular}{lcc}
        \toprule
        \textbf{Component} & \textbf{Dim.} & \textbf{Description} \\
        \midrule
        Hero state         & 4  & HP ratio, block ratio, mana ratio, gold ratio \\
        Active effects     & 7  & Strengh, Weakness, Poison, Burn, Vulnerability, Thorns, Metallicize \\
        Hand cards         & 70 & 10 slots $\times$ 7 features (type, star, cost, val., target, exhaust, affordable) \\
        Deck star counts   & 5  & Count of 1-star through 5-star cards in deck \\
        Enemies            & 24 & 4 slots $\times$ 6 features (HP ratio, block ratio, 4 intent features) \\
        Relics             & 10 & 10 slots: equipped relic IDs (normalized) \\
        Map context        & 7  & Current floor/15, HP ratio, act, shop available, boss type, two flags \\
        Next nodes         & 14 & 7 nodes $\times$ 2 features (room type, estimated difficulty) \\
        \midrule
        \textbf{New features (v2)} & & \\
        Enemy intent damage & 1  & Predicted total incoming damage this turn (normalized) \\
        Affordable hand dmg & 1  & Maximum damage achievable with current mana \\
        Combat turn counter & 1  & Current turn within combat (normalized) \\
        Deck attack ratio   & 1  & Fraction of attack cards in full deck \\
        \midrule
        \textbf{Total} & \textbf{148} & \\
        \bottomrule
    \end{tabular}
\end{table}

All values are normalized to $[0, 1]$ to stabilize neural network training. Unknown or empty slots are filled with zeros (e.g., empty hand slots, absent enemies).

\subsection{Action Space}

The action space is discrete with 68 actions:
\begin{itemize}
    \item Actions 0--9: Play hand-card slot 0--9 (targeting the priority enemy).
    \item Actions 10--49: Play hand-card slot 0--9 targeting enemy slot 0--3 (40 targeted actions for multi-enemy rooms).
    \item Action 50: End player turn (pass to enemy phase).
    \item Actions 51--65: Map navigation -- choose from 15 possible next-room node indices.
    \item Action 66: Use heal item (health potion) acquired in combat.
    \item Action 67: Card reward selection -- choose card index 0 (skip used as action 50 in reward context).
\end{itemize}

At each step, the game engine returns an \textbf{action mask} -- a binary vector $\mathbf{m} \in \{0,1\}^{68}$ -- indicating which actions are currently legal. The mask is computed by querying the C\# game engine for the current valid-action set and converting to Python. Illegal actions (e.g., playing a card with insufficient mana, navigating to an unavailable room) always have mask value 0.

\subsection{Episode Structure}

One \textit{episode} corresponds to one complete game run from Floor 1 to either Floor 15 victory or death. The environment \texttt{reset()} starts a new run (optionally applying Domain Randomization), and \texttt{step(action)} advances the game state by one decision, returning $(s', r, \text{done}, \text{info})$.

\section{RL Agent Persona Design}

\subsection{Design Rationale}

Using a single agent persona for game evaluation has a significant limitation: it may only reveal balance problems that matter to one specific play style. A defensive specialist might expose different vulnerabilities than an aggressive burn-damage specialist. By training four distinct personas, we obtain a richer evaluation signal that better approximates the diversity of real players.

Each persona shares the same observation space (148-dim) and action space (68 actions), but is shaped by a distinct reward function that incentivizes a specific behavioral style. The neural network architecture for all agents is a two-hidden-layer MLP with layers of size $[512, 512]$ for both the actor and critic.

\subsection{Aggressive Agent}

\textbf{Behavioral goal}: Minimize the number of turns to defeat enemies. Accept HP loss in exchange for faster kills. Prefer high-damage cards over defensive ones.

\textbf{Reward function}:
\begin{align}
r_t^{\text{agg}} &= r_{\text{base}} - 0.10 \cdot \mathbb{1}[\text{step taken}] - 0.10 \cdot \Delta\text{HP}_{\text{lost}} \cdot \mathbb{1}[\text{HP ratio} < 0.25] \nonumber \\
&\quad + 5.0 \cdot \mathbb{1}[\text{enemy killed}] + 30.0 \cdot \mathbb{1}[\text{floor cleared}] + 50.0 \cdot \mathbb{1}[\text{victory}] - 100.0 \cdot \mathbb{1}[\text{death}]
\end{align}
The step penalty $-0.10$ drives the agent to end combats quickly. The HP loss penalty is amplified $3\times$ below 25\% HP (a ``non-linear brake'') to prevent suicidal behavior while still permitting calculated risk-taking at full HP.

\subsection{Defensive Agent}

\textbf{Behavioral goal}: Survive as long as possible, prioritizing healing and block generation. Sacrifice offensive tempo to preserve HP.

\textbf{Reward function}:
\begin{align}
r_t^{\text{def}} &= r_{\text{base}} - 0.005 \cdot \mathbb{1}[\text{step taken}] + 0.30 \cdot \Delta\text{HP}_{\text{healed}} - 0.15 \cdot \min(\Delta\text{HP}_{\text{lost}}, 20) \nonumber \\
&\quad + 0.30 \cdot \Delta\text{Block} + 2.0 \cdot \mathbb{1}[\text{combat survived}] + 500.0 \cdot \mathbb{1}[\text{victory}] - 50.0 \cdot \mathbb{1}[\text{death}]
\end{align}
The asymmetric HP weights ($+0.30$ for healing vs.\ $-0.15$ for damage) create a strong preference for rest sites and block cards. The HP-loss cap at 20 per step prevents gradient instability when the agent receives a large boss hit. The very small step penalty ($-0.005$) means the agent is not penalized for taking many turns to kill enemies.

\subsection{Balanced Agent}

\textbf{Behavioral goal}: Represent a standard player who weighs offense and defense equally, manages gold efficiently, and aims for a moderate win rate.

\textbf{Reward function}:
\begin{align}
r_t^{\text{bal}} &= r_{\text{base}} + 0.20 \cdot \Delta\text{HP}_{\text{healed}} - 0.20 \cdot \Delta\text{HP}_{\text{lost}} \nonumber \\
&\quad + 0.02 \cdot \Delta\text{gold} + 30 \cdot \left(\frac{\text{HP}_{\text{final}}}{\text{HP}_{\text{max}}}\right)^2 \cdot \mathbb{1}[\text{victory}] - 50.0 \cdot \mathbb{1}[\text{death}]
\end{align}
The symmetric HP weights model a player who is neither reckless nor overly cautious. The quadratic terminal HP bonus incentivizes the agent to win with as much HP remaining as possible, serving the balancing objective of ``victory HP.''

\subsection{Adaptive Agent}

\textbf{Behavioral goal}: Dynamically switch between aggressive and defensive behaviors based on current HP, without being explicitly programmed to do so.

\textbf{Design challenge}: An earlier version of this agent used a reward function that explicitly changed its formula based on HP ratio at runtime (high HP $\to$ aggressive formula; low HP $\to$ defensive formula). This violates the \textbf{stationarity assumption} of PPO's MDP formulation: the reward function must not depend on hidden state outside the observation. Non-stationary rewards cause the value function to fail to converge, as the same $(s, a)$ pair can yield different expected future rewards depending on unobservable context.

\textbf{Revised design (v4)}: The Adaptive Agent uses a minimal, stationary reward function:
\begin{equation}
r_t^{\text{ada}} = r_{\text{base}} + 3.0 \cdot \mathbb{1}[\text{enemy killed}] - 0.10 \cdot \Delta\text{HP}_{\text{lost}} - 0.02 \cdot \mathbb{1}[\text{step taken}]
\end{equation}
Crucially, the HP ratio is included as the \textbf{first element of the observation vector} $s[0]$. With a large enough network ($[512, 512]$), the agent learns an \textbf{emergent strategy-switching} behavior: it discovers independently that attacking aggressively when $s[0] > 0.7$ and switching to block-heavy defense when $s[0] < 0.3$ maximizes the stationary reward. This \textit{emergent} adaptation is more robust than hand-coded switching because the network can condition on the full observation, not just HP.

\section{Optimization Layer Design}

\subsection{Data Flow}

Figure~\ref{fig:dataflow} shows how data flows from the optimization layer through simulation to fitness evaluation.

\begin{figure}[H]
    \centering
    \includegraphics[width=0.9\textwidth]{dataflow.png}
    \caption{Data flow in the optimization loop. A genome is created by the optimization layer, applied to a scaled \texttt{GameData} copy, passed to the simulation runner which uses the AI agent to play games, and resulting statistics are evaluated into fitness scores that feed back to the optimizer.}
    \label{fig:dataflow}
\end{figure}

\subsection{Genome Representations}

\subsubsection{BalanceGenome (Single-Objective GA)}

The \texttt{BalanceGenome} directly encodes one scalar multiplier per game parameter, resulting in a chromosome of approximately 500 dimensions. The genome consists of four dictionaries:
\begin{itemize}
    \item \texttt{CardActionScalars}: One float per card action value (damage or block amount per card).
    \item \texttt{CardCostModifiers}: Integer adjustments to each card's mana cost.
    \item \texttt{EnemyHealthScalars}: One float per enemy type scaling base HP.
    \item \texttt{EnemyActionValueScalars}: One float per enemy action scaling attack damage.
\end{itemize}
This flat representation allows fine-grained per-entity adjustment but suffers from the curse of dimensionality and tends to produce incoherent solutions where logically related parameters (e.g., all 1-star cards) receive unrelated scalars.

\subsubsection{HierarchicalGenome (NSGA-II)}

The \texttt{HierarchicalGenome} reduces the search space to approximately 50 dimensions by grouping parameters into semantically meaningful categories:
\begin{itemize}
    \item \textbf{Global multipliers} (5 params): Overall scaling for damage, health, block, mana cost, gold rewards.
    \item \textbf{Progression scaling} (9 params): Separate damage, health, and block multipliers for Early (floors 1--5), Mid (floors 6--10), and Late (floors 11--15) game phases.
    \item \textbf{Category scaling} (8 params): Per-type scaling for Attack/Skill/Power cards; per-star-rating scaling for both cards and enemies.
    \item \textbf{Room distribution} (5 params): Relative weights for Monster, Elite, Event, Shop, and Rest room types.
    \item \textbf{Miscellaneous} (8 params): Hero health scalar, rest-site healing rate, start gold, per-star enemy ratios, difficulty progression rate.
\end{itemize}
The hierarchical structure ensures that NSGA-II can efficiently navigate the search space by modifying high-level properties (``make late-game enemies 20\% tougher'') rather than individually tuning hundreds of entity-level values.

\subsection{Fitness Functions}

The same three sub-objectives are computed for both GA (aggregated) and NSGA-II (separately):

\textbf{Balance Score} uses Gaussian scoring to evaluate how close the simulation metrics are to the target values:
\begin{equation}
f_{\text{balance}} = \frac{1}{3}\left(e^{-\alpha_{wr} \cdot \Delta_{wr}^2} + e^{-\alpha_{hp} \cdot \Delta_{hp}^2} + e^{-\alpha_{fd} \cdot \Delta_{fd}^2}\right)
\label{eq:balance_score}
\end{equation}
where $\Delta_{wr} = |\text{WinRate} - 0.45|$, $\Delta_{hp} = |\text{VictoryHP} - 0.30|$, $\Delta_{fd} = |\text{FloorOnDeath} - 10|$, and $\alpha_{wr}=10.0$, $\alpha_{hp}=8.0$, $\alpha_{fd}=0.05$.

\textbf{Engagement Score} rewards card pool diversity:
\begin{equation}
f_{\text{engagement}} = \alpha \cdot \frac{|\text{ViableCards}|}{|\text{TotalCards}|} + \beta \cdot \text{BuildDiversity}
\label{eq:engagement_score}
\end{equation}
A card is ``viable'' if its pick rate exceeds a threshold $\tau$ and its conditional win rate is above a second threshold $\tau_w$. \texttt{BuildDiversity} measures the entropy of the card-selection distribution.

\textbf{Coherence Score} penalizes trap cards and extreme parameter values:
\begin{equation}
f_{\text{coherence}} = 1 - \text{TrapCardPenalty} - \lambda \sum_{j}\left(\text{scalar}_j - 1\right)^2
\label{eq:coherence_score}
\end{equation}
A trap card is a frequently picked card whose conditional win rate is significantly \textit{below} the population average. The regularization term $\lambda \sum_j (\text{scalar}_j - 1)^2$ penalizes solutions that deviate far from the baseline, preventing extreme parameter configurations.
